\documentclass[11pt]{article}
\usepackage[T1]{fontenc}
\usepackage[utf8]{inputenc}
\usepackage{lmodern,amsmath,amssymb,amsthm,mathtools,booktabs,microtype}
\usepackage[margin=1in]{geometry}
\usepackage[colorlinks=true,linkcolor=blue,citecolor=blue,urlcolor=blue]{hyperref}
\hypersetup{pdftitle={At least one of beta(4), beta(6), ..., beta(32) is irrational},
 pdfauthor={twinheterodox},pdfsubject={An isolated-pole denominator refinement for even Dirichlet beta values}}
\newtheorem{theorem}{Theorem}[section]
\newtheorem{proposition}[theorem]{Proposition}
\newtheorem{lemma}[theorem]{Lemma}
\theoremstyle{definition}
\newtheorem{example}[theorem]{Example}
\newcommand{\Q}{\mathbb Q}
\newcommand{\Z}{\mathbb Z}
\newcommand{\R}{\mathbb R}
\newcommand{\vp}{v_p}
\newcommand{\ind}{\mathbf 1}
\newcommand{\ssum}{\mathop{\sum}}
\numberwithin{equation}{section}
\title{At least one of $\beta(4),\beta(6),\ldots,\beta(32)$ is irrational}
\author{twinheterodox\\[2pt]\normalsize Prepared with Codex}
\date{11 November 2026}
\begin{document}
\maketitle
\begin{abstract}
For the Dirichlet beta function, we prove that at least one of
$\beta(4),\beta(6),\ldots,\beta(32)$ is irrational. This reduces
Krisman's seventeen-value set $\beta(4),\beta(6),\ldots,\beta(36)$
\cite{krisman} to fifteen values. The improvement comes from a sharper
denominator estimate for his construction.
\end{abstract}
\input{paper/sections/01-introduction}
\input{paper/sections/02-forms}
\input{paper/sections/03-denominators}
\input{paper/sections/04-decay}
\input{paper/sections/05-nonvanishing}
\input{paper/sections/06-computations}
\appendix
\input{paper/sections/appendix-small-primes}
\input{paper/sections/references}
\end{document}
